% Source: public/notes/物工/問題/2026ap_exam.pdf, Japanese PDF pp. 4--5.
\documentclass[11pt,a4paper]{article}
\usepackage[margin=24mm]{geometry}
\usepackage{fontspec}
\usepackage{xeCJK}
\setCJKmainfont{HaranoAjiMincho-Regular.otf}
\setCJKsansfont{HaranoAjiGothic-Medium.otf}
\usepackage{amsmath,amssymb}
\setlength{\parindent}{1em}
\setlength{\parskip}{0.35em}
\title{物理工学専攻 院試解答集\\2026年 物理学 第2問}
\author{政岡凜太郎}
\date{}
\begin{document}
\maketitle

\section*{[1.1] 自由回転状態の規格化}
角度 $\phi$ は一周すると同じ配置に戻る。従って一周分 $0\leq\phi<2\pi$ で規格化する。$A$ は正の実数だから
\begin{equation}
1=\int_0^{2\pi}|\Psi_m(\phi)|^2\,\mathrm d\phi
 =A^2\int_0^{2\pi}\mathrm d\phi=2\pi A^2,
\qquad \boxed{A=\frac{1}{\sqrt{2\pi}}}.
\end{equation}

\section*{[1.2] 周期性とエネルギー}
波動関数の一価性 $\Psi_m(\phi+2\pi)=\Psi_m(\phi)$ は $e^{2\pi i m}=1$、すなわち $m\in\mathbb Z$ を要求する。$V_3=0$ の (2) に $e^{im\phi}$ を代入すると、二階微分は $-m^2e^{im\phi}$ なので
\begin{equation}
\boxed{E_m=\frac{\hbar^2m^2}{2I}},\qquad m=0,\pm1,\pm2,\ldots.
\end{equation}
$m\neq0$ では $m$ と $-m$ が同じエネルギーを持つ。

\section*{[2.1] 極小近傍の調和振動}
極小は $\phi_k=2\pi k/3$ ($k=0,1,2$) であり、$V(\phi_k)=0$ である。$\phi-\phi_k$ が小さいとき $1-\cos[3(\phi-\phi_k)]\simeq\tfrac92(\phi-\phi_k)^2$ だから
\begin{equation}
V(\phi)\simeq\frac{9V_3}{4}(\phi-\phi_k)^2
=\frac12 I\omega_0^2(\phi-\phi_k)^2.
\end{equation}
回転の運動エネルギーは $I\dot\phi^2/2$ なので、調和振動子の質量に相当する係数は $I$ である。従って
\begin{equation}
\boxed{\omega_0=3\sqrt{\frac{V_3}{2I}}},\qquad
\boxed{\epsilon=\frac{\hbar\omega_0}{2}=\frac{3\hbar}{2}\sqrt{\frac{V_3}{2I}}}.
\end{equation}

\section*{[2.2] 三つの極小点を結ぶトンネル効果}
以下、(3) と (6) の通常の行列固有値問題として、局在状態 $|1\rangle,|2\rangle,|3\rangle$ を互いに直交する規格化基底として扱う。原題は各状態の規格化だけを明記しているが、重なりが無視できない場合には別途重なり行列を要する。

(3) の各成分は
\begin{equation}
(\widetilde H\boldsymbol c)_j=(\epsilon+t)c_j-t(c_1+c_2+c_3),
\qquad \boldsymbol c=\begin{pmatrix}c_1\\c_2\\c_3\end{pmatrix}
\end{equation}
と書ける。三成分が等しい対称状態では $c_1=c_2=c_3=e^{i\alpha}/\sqrt3$、$\alpha$ は任意の全体位相であり、固有値は $\boxed{E_{\mathrm s}=\epsilon-2t}$ となる。一方、$c_1+c_2+c_3=0$ の二次元部分空間では $\boxed{E_{\mathrm d}=\epsilon+t}$ が二重縮退する。この部分空間の任意の状態は $|c_1|^2+|c_2|^2+|c_3|^2=1$ を満たせば規格化される。三つの係数を個別に一意に定める条件はない。

\section*{[2.3] 回転操作との同時固有状態}
(4) から、$R$ は係数を $(c_1,c_2,c_3)\mapsto(c_3,c_1,c_2)$ と移す。$\omega^3=1$ および $1+\omega+\omega^2=0$ を用いれば、規格化された同時固有状態を
\begin{align}
\boxed{|\Psi_1\rangle}&=\frac{|1\rangle+|2\rangle+|3\rangle}{\sqrt3},
& R|\Psi_1\rangle&=|\Psi_1\rangle,
& \widetilde H|\Psi_1\rangle&=(\epsilon-2t)|\Psi_1\rangle,\\
\boxed{|\Psi_{\omega}\rangle}&=\frac{|1\rangle+\omega^2|2\rangle+\omega|3\rangle}{\sqrt3},
& R|\Psi_{\omega}\rangle&=\omega|\Psi_{\omega}\rangle,
& \widetilde H|\Psi_{\omega}\rangle&=(\epsilon+t)|\Psi_{\omega}\rangle,\\
\boxed{|\Psi_{\omega^2}\rangle}&=\frac{|1\rangle+\omega|2\rangle+\omega^2|3\rangle}{\sqrt3},
& R|\Psi_{\omega^2}\rangle&=\omega^2|\Psi_{\omega^2}\rangle,
& \widetilde H|\Psi_{\omega^2}\rangle&=(\epsilon+t)|\Psi_{\omega^2}\rangle
\end{align}
と選べる。例えば $R|\Psi_{\omega}\rangle=(\omega|1\rangle+|2\rangle+\omega^2|3\rangle)/\sqrt3=\omega|\Psi_{\omega}\rangle$ であり、位相の対応を直接確かめられる。三状態は互いに直交する。

\section*{[2.4] WKB の障壁積分}
$\epsilon\simeq0$ とすると、隣接する極小間の転回点は $\phi_1\simeq0$、$\phi_2\simeq2\pi/3$ となる。この区間では $\sin(3\phi/2)\geq0$ であり、$1-\cos(3\phi)=2\sin^2(3\phi/2)$ だから、(5) の作用積分は
\begin{align}
\int_0^{2\pi/3}\!\sqrt{2I V(\phi)}\,\mathrm d\phi
&=\sqrt{2IV_3}\int_0^{2\pi/3}\sin\!\left(\frac{3\phi}{2}\right)\mathrm d\phi\\
&=\sqrt{2IV_3}\left[-\frac23\cos\!\left(\frac{3\phi}{2}\right)\right]_0^{2\pi/3}
=\frac43\sqrt{2IV_3}.
\end{align}
従って深い障壁 $V_3\gg\epsilon$、すなわち $\sqrt{2IV_3}\gg\hbar$ の指数近似で
\begin{equation}
\boxed{t\approx C\exp\!\left[-\frac{4\sqrt{2IV_3}}{3\hbar}\right]}.
\end{equation}
$\sqrt{IV_3}$ は作用の次元を持つので指数は無次元である。$C$ は (5) で未指定の前因子であり、この近似だけでは決められない。

\section*{[3.1] 対称・反対称基底での行列}
$|A\rangle=(|2\rangle+|3\rangle)/\sqrt2$、$|B\rangle=(|2\rangle-|3\rangle)/\sqrt2$ とする。$|1\rangle$ は $|2\rangle,|3\rangle$ に同じ $-t'$ で結合するので、$|A\rangle$ への行列要素は $-\sqrt2t'$、$|B\rangle$ への行列要素は零である。また $|2\rangle,|3\rangle$ 間の $-t$ は対称・反対称状態をそれぞれ $\epsilon-t,\epsilon+t$ に分ける。従って指定された順序 $(|1\rangle,|A\rangle,|B\rangle)$ で
\begin{equation}
\boxed{\widetilde H'=
\begin{pmatrix}
\epsilon-\Delta&-\sqrt2t'&0\\
-\sqrt2t'&\epsilon-t&0\\
0&0&\epsilon+t
\end{pmatrix}}.
\end{equation}

\section*{[3.2] 非対角項の二次摂動}
指定の分割は
\begin{equation}
\widetilde H_0=\operatorname{diag}(\epsilon-\Delta,\epsilon-t,\epsilon+t),\qquad
\widetilde H_1=\begin{pmatrix}0&-\sqrt2t'&0\\-\sqrt2t'&0&0\\0&0&0\end{pmatrix}.
\end{equation}
$\widetilde H_1$ の対角要素は零なので一次補正は零である。$\Delta>t$ により、非摂動の $|1\rangle$ と $|A\rangle$ のエネルギー差は $(\epsilon-\Delta)-(\epsilon-t)=-(\Delta-t)$ と零でない。以下は$t'$についての形式的な二次展開であり、近似値として用いる条件は$\sqrt2t'\ll\Delta-t$である。二次の式 $E_n^{(2)}=\sum_{k\ne n}|\langle k|\widetilde H_1|n\rangle|^2/(E_n^{(0)}-E_k^{(0)})$ を用いると
\begin{align}
\boxed{E_{1}}&=\epsilon-\Delta-\frac{2t'^2}{\Delta-t}+O\!\left(\frac{t'^4}{(\Delta-t)^3}\right),\\
\boxed{E_{A}}&=\epsilon-t+\frac{2t'^2}{\Delta-t}+O\!\left(\frac{t'^4}{(\Delta-t)^3}\right),\\
\boxed{E_{B}}&=\epsilon+t.
\end{align}
$|B\rangle$ は他の状態と結合せず、最後の値はこの行列内で厳密である。二次展開の条件は $\sqrt2t'\ll\Delta-t$ であり、原題の大小関係 $\Delta>t>t'>0$ だけでは $\Delta-t$ が十分大きいとは限らない。実際、最初の二準位の厳密値は
\begin{equation}
E_{\pm}=\epsilon-\frac{\Delta+t}{2}
\pm\sqrt{\frac{(\Delta-t)^2}{4}+2t'^2},
\end{equation}
であり、上の $E_1,E_A$ はそれぞれ $E_-,E_+$ の小さい $t'/(\Delta-t)$ に関する展開である。
\end{document}
